\section{Sujet n\textsuperscript{o}\,8}
\noindent\begin{minipage}{0.5\linewidth}\footnotesize Office du Baccalauréat du Cameroun\end{minipage}%
\hfill\begin{minipage}{0.45\linewidth}\footnotesize\raggedleft Examen : \textbf{Baccalauréat} · Série : \textbf{C/E}\\
Épreuve : \textbf{Mathématiques} · Durée : \textbf{4\,h} · Coef. : \textbf{7}\end{minipage}
\par\smallskip{\color{onbo}\rule{\linewidth}{1pt}}

\subsection*{Partie A — Évaluation des ressources \hfill (15 points)}

\noindent\textbf{Exercice 1.} \hfill\textit{(3 points)}\\
Un pisciculteur de Foumban dispose de $3$ bassins. Le bassin $B_1$ contient $5$ tilapias
et $3$ silures, le bassin $B_2$ contient $4$ tilapias et $4$ silures. On choisit un bassin
au hasard ($B_1$ ou $B_2$, équiprobables), puis on y pêche simultanément $2$ poissons.
\begin{enumerate}[label=\arabic*)]
\item Calculer la probabilité de pêcher $2$ tilapias dans $B_1$, puis dans $B_2$. \hfill\textit{(1 pt)}
\item En déduire la probabilité $p$ de pêcher $2$ tilapias. \hfill\textit{(1 pt)}
\item Sachant qu'on a pêché $2$ tilapias, calculer la probabilité que ce soit dans $B_1$. \hfill\textit{(1 pt)}
\end{enumerate}

\smallskip
\noindent\textbf{Exercice 2.} \hfill\textit{(3 points)}\\
Pour numéroter ses alevins, le pisciculteur utilise un code à deux entiers $(a\,;\,b)$
vérifiant $\operatorname{PGCD}(a,b)=12$ et $\operatorname{PPCM}(a,b)=180$.
\begin{enumerate}[label=\arabic*)]
\item Montrer que tout couple solution s'écrit $a=12a'$, $b=12b'$ avec $a'\wedge b'=1$ et $a'b'=15$. \hfill\textit{(1,5 pt)}
\item En déduire tous les couples $(a\,;\,b)$ tels que $a\le b$. \hfill\textit{(1,5 pt)}
\end{enumerate}

\smallskip
\noindent\textbf{Exercice 3.} \hfill\textit{(4 points)}\\
Le taux d'oxygène (en mg/L) d'un bassin évolue selon $f(t)=4+t\,\mathrm{e}^{-t/2}$ où $t\ge0$
est le temps en heures, de courbe $(\mathcal{C})$.
\begin{enumerate}[label=\arabic*)]
\item Déterminer la limite de $f$ en $+\infty$ ; interpréter graphiquement. \hfill\textit{(1 pt)}
\item Montrer que $f'(t)=\left(1-\tfrac{t}{2}\right)\mathrm{e}^{-t/2}$ et dresser le tableau de variations de $f$ sur $[0\,;+\infty[$. \hfill\textit{(1,5 pt)}
\item À l'aide d'une intégration par parties, montrer que $\displaystyle\int_0^{2}t\,\mathrm{e}^{-t/2}\,\dd t=4-8\mathrm{e}^{-1}$. \hfill\textit{(1,5 pt)}
\end{enumerate}

\smallskip
\noindent\textbf{Exercice 4.} \hfill\textit{(5 points)}\\
Le plan complexe est muni d'un repère orthonormé direct $(O\,;\,\vec u,\vec v)$. On considère
les points $A$, $B$, $C$ d'affixes $z_A=2$, $z_B=2i$, $z_C=-1+i$.
\begin{enumerate}[label=\arabic*)]
\item Calculer $\dfrac{z_C-z_A}{z_B-z_A}$ sous forme algébrique. \hfill\textit{(1 pt)}
\item En déduire la nature du triangle $ABC$ et une mesure de l'angle $\widehat{BAC}$. \hfill\textit{(1,5 pt)}
\item Soit $S$ la similitude directe de centre $A$, d'angle $\frac\pi2$ et de rapport $\sqrt2$.
Donner l'écriture complexe de $S$. \hfill\textit{(1,5 pt)}
\item Déterminer l'affixe de l'image $B'$ de $B$ par $S$. \hfill\textit{(1 pt)}
\end{enumerate}

\subsection*{Partie B — Évaluation des compétences \hfill (5 points)}

\noindent\textit{Monsieur Njoya, pêcheur et pisciculteur au lac de Bamendjing, sollicite ton
aide, en tant qu'élève de Terminale~C, pour trois décisions. Il veut d'abord construire un
bassin rectangulaire dont l'un des grands côtés s'appuie sur une digue déjà bétonnée ; il
dispose de \SI{60}{\metre} de margelle pour les trois autres côtés. Sa première récolte
annuelle est de \SI{800}{\kilogram} de poisson ; grâce à un meilleur alevinage, elle augmente
de \SI{15}{\percent} chaque année. Enfin, à la livraison, $6\,\%$ des poissons arrivent
abîmés ; un grossiste en contrôle $5$ pris au hasard.}

\begin{enumerate}[label=\textbf{Tâche \arabic*.},leftmargin=2.4em]
\item Déterminer les dimensions du bassin rectangulaire d'\emph{aire maximale} qu'il peut
construire, ainsi que cette aire. \hfill\textit{(1,5 pt)}
\item Déterminer l'année à partir de laquelle sa récolte annuelle dépassera \SI{1500}{\kilogram}. \hfill\textit{(1,5 pt)}
\item Déterminer la probabilité qu'\emph{au moins un} des $5$ poissons contrôlés soit abîmé. \hfill\textit{(1,5 pt)}
\end{enumerate}
\par\smallskip{\footnotesize\itshape Présentation générale : 0,5 point.}

% ════════════════════════════════════════════════════
\corrige{du sujet 8}

\noindent\textbf{Partie A — Exercice 1.}
1) Dans $B_1$ : $\binom82=28$ tirages, $P_1=\dfrac{\binom52}{28}=\dfrac{10}{28}=\dfrac{5}{14}$.
Dans $B_2$ : $P_2=\dfrac{\binom42}{28}=\dfrac{6}{28}=\dfrac{3}{14}$.
2) Bassin équiprobable : $p=\frac12 P_1+\frac12 P_2=\frac12\big(\frac5{14}+\frac3{14}\big)=\frac12\cdot\frac8{14}=\dfrac{2}{7}$.
3) Bayes : $P(B_1\mid 2T)=\dfrac{\frac12 P_1}{p}=\dfrac{\frac12\cdot\frac5{14}}{\frac27}=\dfrac{5/28}{8/28}=\dfrac58$.

\noindent\textbf{Exercice 2.}
1) Si $d=\operatorname{PGCD}(a,b)=12$, on écrit $a=12a'$, $b=12b'$ avec $a'\wedge b'=1$.
Or $\operatorname{PGCD}\cdot\operatorname{PPCM}=ab$, donc $12\times180=144\,a'b'$, soit $a'b'=\frac{2160}{144}=15$.
2) Couples $(a',b')$ premiers entre eux, produit $15$, avec $a'\le b'$ : $(1\,;15)$ et $(3\,;5)$
(le couple $(1,15)$ : $1\wedge15=1$ ; $(3,5)$ : $3\wedge5=1$). D'où
$(a\,;b)=(12\,;180)$ et $(36\,;60)$.

\noindent\textbf{Exercice 3.}
1) En $+\infty$, par croissances comparées $t\,\mathrm{e}^{-t/2}\to0$, donc $f(t)\to4$ : la droite
$y=4$ est asymptote horizontale.
2) $f'(t)=\mathrm{e}^{-t/2}+t\cdot(-\tfrac12)\mathrm{e}^{-t/2}=\big(1-\tfrac{t}{2}\big)\mathrm{e}^{-t/2}$.
$f'>0$ pour $t<2$, $f'<0$ pour $t>2$ : $f$ croît sur $[0\,;2]$, décroît sur $[2\,;+\infty[$,
maximum $f(2)=4+2\mathrm{e}^{-1}\approx4{,}74$.
3) IPP ($u=t$, $u'=1$, $v'=\mathrm{e}^{-t/2}$, $v=-2\mathrm{e}^{-t/2}$) :
$\int_0^2 t\,\mathrm{e}^{-t/2}\dd t=\big[-2t\mathrm{e}^{-t/2}\big]_0^2+\int_0^2 2\mathrm{e}^{-t/2}\dd t
=-4\mathrm{e}^{-1}+\big[-4\mathrm{e}^{-t/2}\big]_0^2=-4\mathrm{e}^{-1}+(-4\mathrm{e}^{-1}+4)=4-8\mathrm{e}^{-1}$.

\noindent\textbf{Exercice 4.}
1) $\dfrac{z_C-z_A}{z_B-z_A}=\dfrac{-1+i-2}{2i-2}=\dfrac{-3+i}{-2+2i}$.
On multiplie par le conjugué : $\dfrac{(-3+i)(-2-2i)}{(-2+2i)(-2-2i)}=\dfrac{6+6i-2i-2i^2}{4+4}=\dfrac{6+4i+2}{8}=\dfrac{8+4i}{8}=1+\tfrac12 i$.
2) Le quotient vaut $1+\frac12 i$, de module $\sqrt{1+\frac14}=\frac{\sqrt5}{2}\ne1$ et d'argument
$\theta=\arctan\frac12$. Comme l'argument n'est pas $\pm\frac\pi2$, le triangle n'est pas rectangle
en $A$ de cette manière ; il est quelconque, et $\widehat{BAC}=\arctan\tfrac12\approx26{,}6^\circ$.
(On a $|z_B-z_A|=2\sqrt2$ et $|z_C-z_A|=\sqrt{10}$, donc $AC=AB\cdot\frac{\sqrt5}{2}$.)
3) $S(z)=a z+b$ avec $a=\sqrt2\,\mathrm{e}^{i\pi/2}=\sqrt2\,i$ et centre $A$ fixe : $z_A=a z_A+b$, donc
$b=z_A(1-a)=2(1-\sqrt2 i)=2-2\sqrt2 i$. Ainsi $S(z)=\sqrt2 i\,z+2-2\sqrt2 i$.
4) $z_{B'}=\sqrt2 i\cdot 2i+2-2\sqrt2 i=2\sqrt2 i^2+2-2\sqrt2 i=-2\sqrt2+2-2\sqrt2 i
=(2-2\sqrt2)-2\sqrt2\, i$.

\smallskip
\noindent\textbf{Partie B — Situation.}
\emph{Tâche 1.} Le grand côté $\ell$ s'appuie sur la digue : on clôture deux largeurs $x$ et une
longueur $\ell$, donc $2x+\ell=60$, soit $\ell=60-2x$. Aire $A(x)=x(60-2x)$ ;
$A'(x)=60-4x=0$ en $x=15$. Dimensions : largeur $\SI{15}{\metre}$, longueur $\SI{30}{\metre}$ ;
aire maximale $A=15\times30=\SI{450}{\metre\squared}$.

\emph{Tâche 2.} Récolte de l'année $n$ : $u_n=800\times1{,}15^{\,n-1}$. On résout $u_n>1500$,
soit $1{,}15^{\,n-1}>\frac{1500}{800}=1{,}875$, d'où $n-1>\dfrac{\ln1{,}875}{\ln1{,}15}\approx\dfrac{0{,}6286}{0{,}1398}\approx4{,}5$ :
à partir de la \textbf{6\textsuperscript{e} année} ($u_6=800\times1{,}15^5\approx\SI{1609}{\kilogram}$).

\emph{Tâche 3.} $P(\text{au moins un abîmé})=1-(0{,}94)^5\approx1-0{,}734=\mathbf{0{,}266}$, soit environ \SI{26,6}{\percent}.
